\begin{abstract}
We prove that the complex border determinantal complexity of the $m\times m$
permanent is $\Omega(m^3)$, allowing arbitrary affine-linear determinant
representations and coefficientwise limits. It also gives cubic lower bounds
for exact determinantal complexity and for the numbers of vertices and edges
in affine-linear algebraic branching programs, including coefficientwise
limits with a fixed vertex or edge budget.
\end{abstract}

\section{Introduction}
\label{sec:introduction}

For a matrix $X=(x_{ij})$ of independent variables, the permanent is
\[
 \per_m(X)=\sum_{\sigma\in S_m}\prod_{i=1}^m x_{i,\sigma(i)}.
\]
The permanent--determinant problem asks how large a determinant of
affine-linear forms is needed to express $\per_m$. Valiant connected
determinant projections to arithmetic formulas: his construction represents
a polynomial computed by a formula with $s$ addition or multiplication
operations as a determinant of order $s+2$
\cite[Theorem~1]{Valiant1979Algebra}. His algebraic completeness theorem
shows that, in characteristic different from two, polynomial families
defined by summing polynomial-size formulas over Boolean strings of
polynomial length are projections of polynomial-size permanents
\cite[Theorem~2]{Valiant1979Algebra}. Valiant's projections substitute
individual variables or field constants. The exact model below allows
every matrix entry to be an arbitrary affine-linear form.

For $f\in\C[x_1,\ldots,x_N]$, its \emph{determinantal complexity}
$\dc(f)$ is the least positive integer $n$ for which
\[
 f=\det\left(A_0+\sum_{j=1}^N x_jA_j\right),
 \qquad A_0,\ldots,A_N\in\Mat_n(\C).
\]
Its \emph{border determinantal complexity} $\bdc(f)$ is the least positive
integer $n$ such that $f$ lies in the coefficientwise closure of
\begin{equation}
 \left\{\det\left(A_0+\sum_{j=1}^N x_jA_j\right):
 A_0,\ldots,A_N\in\Mat_n(\C)\right\}.
 \label{eq:border-definition}
\end{equation}
The closure is taken in the finite-dimensional space of polynomials of
degree at most $n$; a polynomial of larger degree cannot belong to it.
We allow $f$ to be constant or zero. Coefficientwise closure agrees with
Zariski closure here, as explained in Section~\ref{sec:models}. The size
$n$ stays fixed along an approximating sequence, while the complex
coefficients of its matrix entries are unrestricted.

\begin{theorem}
\label{thm:main}
There exist absolute constants $c>0$ and $m_0$ such that, for every integer
$m\ge m_0$,
\[
 \bdc(\per_m)\ge c m^3.
\]
One may take $c=1/(5\,529\,600e)$ and $m_0=1408$.
\end{theorem}

The constants are not optimized. Proposition~\ref{prop:model-equivalence}
identifies this coefficientwise model with the projective determinant-orbit
closure formulation, including equal permanent and determinant orders.

The same numerical lower bound holds for exact determinantal complexity.
There is also a consequence for algebraic branching programs. An
affine-linear algebraic branching program is a finite directed acyclic
graph with a source and a sink and an affine-linear form on each edge; it
computes the sum of products of edge labels along all source-to-sink paths.
Over $\C$, computing $\per_m$ requires $\Omega(m^3)$ vertices and
$\Omega(m^3)$ edges. The same bounds hold for coefficientwise limits when
the relevant vertex or edge budget remains fixed along the approximation,
even if the graph varies. Section~\ref{sec:consequences} gives the precise
statements and determinant reductions.

